Worst-group audits are meant to answer an operational question: is there a subgroup for which a deployed model is unreliable enough to require action? Under covariate shift, the usual answer is to importance-weight labeled source data to represent the target population, then report the subgroup-bin cell with the largest estimated residual. This paper shows that such alarms can be statistically misleading. The relevant question is not only whether the shifted estimand is transportable, but whether the selected cell is locally supported by enough weighted labels to certify the alarm.

We develop a compact inference framework for worst-group reliability under covariate shift. The first result is an identifiability boundary: if a target subgroup-bin cell has no source support, its target reliability cannot be certified from labeled source data and unlabeled target covariates, even when covariate shift is otherwise correct. The main result gives simultaneous confidence certificates for a finite family of subgroup-bin cells. The radius is governed by cellwise effective sample size, not global effective sample size. The same theorem covers cross-fitted learned density ratios by adding an explicit local nuisance term, and separates conditional-on-covariates guarantees from population-target guarantees.

Experiments are chosen to test the statistical claim directly. In a calibrated covariate-shift benchmark, naive worst-cell audits and global-ESS gates produce frequent false alarms, while local simultaneous certificates control them and retain power when a planted subgroup failure has adequate overlap. A learned-weight and adaptive-search study shows how density-ratio error and subgroup search add uncertainty beyond oracle local ESS. Two real-data deployment-style audits, Adult Census and ACSIncome, show the practical consequence: large naive subgroup alarms would trigger escalation, but local certificates correctly report insufficient evidence when the alarm-driving cells are thinly supported. The results support a reporting rule: under covariate shift, worst-group alarms should be reported as actionable only when they are locally supported and statistically certified.
